\section{Literature Review}

Research comparing client-side rendering and server-side rendering usually focuses on loading speed, user experience, and search visibility rather than energy use. \citet{nordstrom2023} report that SSR tends to help more as pages become larger and more complex, while the difference is smaller on lighter pages. \citet{clapton2025} describe a similar pattern: SSR performs better on metrics such as TTI and LCP as data volume increases, while CSR can still compare favourably on earlier paint-oriented metrics in lighter cases.

Another recurring theme is that SSR is often valued for what it delivers before the browser has done much work. \citet{ojala2025} describes SSR as beneficial for initial delivery and SEO, while also noting the added complexity that comes with hydration and server involvement. \citet{gaddam2025} similarly presents SSR as useful for SEO and user engagement. Taken together, these studies suggest that render strategy is not only a speed question, but also an architectural trade-off between delivery behaviour, implementation complexity, and where the workload is placed.

Some work moves closer to the concerns of this report by looking at browser-side resource use. \citet{gaddam2025} argues that shifting more work to the server can reduce client-side CPU and memory pressure. \citet{clapton2025} also report heavier JavaScript heap use on the client as dataset size increases, especially when more of the rendering work remains in the browser. This matters because it connects render strategy to the computational load placed on the user's device, not just to visible timing metrics.

Direct energy measurement remains much rarer in SSR/CSR comparisons. \citet{thiagarajan2012} provide an important earlier reference point by showing that JavaScript and CSS processing can account for a substantial part of mobile browser energy use. Their study predates modern React-style rendering patterns, but it is still relevant because it shows that client-side execution itself has a meaningful battery cost.

Recent work also shows that web-energy measurement is methodologically difficult. \citet{kalliola2025} argue that estimating the energy consumption of websites remains challenging and that simple online calculators are often unreliable. Their discussion helps explain why careful measurement design matters when comparing rendering strategies, especially if the goal is to say something more precise than ``this feels faster'' or ``this seems greener''.

Overall, the literature suggests a clear gap. SSR and CSR have been compared extensively as delivery strategies, but much less often with mobile client-side energy use as the main outcome. Resource indicators such as CPU and memory do appear in the literature, yet they are usually secondary to speed and SEO. This study therefore builds on existing performance research while shifting the focus toward a narrower question: how the location of rendering work affects battery-derived energy use and selected browser-side metrics in one shared Next.js prototype.
